\chapter{Rémunération par métier, type d'entreprise et ancienneté}\label{ch:in:pay-levels-by-role-firm-type-and-seniority}

Chaque source publique de rémunération n'en voit qu'une tranche. Les dossiers de visa américains voient les \omterm{def:in:how-pay-works:base}{salaires de base} proposés pour les postes parrainés, et
rien des bonus. L'enquête professionnelle du gouvernement voit les salaires par secteur et par métier, mais exclut la plupart
des bonus et plafonne ce qu'elle publie. Les comptes déposés voient le coût total du personnel par tête, bonus compris, mais font une moyenne qui efface les
métiers. Les rapports de rémunération des banques ne comptent que les personnes payées plus d'un million d'euros. Mises côte à côte, les tranches
divergent d'un facteur dix et plus, et chacune a raison sur ce qu'elle mesure~: en 2025, le \omterm{def:in:how-pay-works:base}{salaire de base} médian proposé
dans les demandes parrainées chez les teneurs de marché de ce livre était de 175\,000~\$, tandis que l'un d'eux dépensait 2.1~millions de dollars par
salarié au Royaume-Uni. Ce chapitre construit des fourchettes à partir de chaque tranche, avec sa mise en garde, par métier, par type d'employeur et par
ancienneté, puis les réunit en un seul chiffre que le lecteur souhaite le plus et qu'aucune source ne publie~: le rapport entre la \omterm{def:in:how-pay-works:base}{rémunération totale} et
le \omterm{def:in:how-pay-works:base}{salaire de base}.

\section{Ce que mesure chaque source}

Cinq sources publiques fournissent des chiffres de rémunération pour le secteur. Chacune mesure une grandeur différente sur une population différente, et
aucune ne mesure ce que le lecteur entend par «~rémunération~».

\begin{definition}[labor condition application, salaire en vigueur (prevailing wage)]\label{def:in:pay-levels-by-role-firm-type-and-seniority:lca}
Une \emph{labor condition application}\index{labor condition application} (LCA) est l'attestation qu'un employeur américain dépose
auprès du Department of Labor avant de parrainer un salarié pour un visa H-1B, H-1B1 ou E-3~: l'intitulé du poste, le code de profession,
le lieu de travail, le salaire proposé ou la fourchette de salaire, et le salaire en vigueur de la profession et de la zone.
Le \emph{salaire en vigueur}\index{salaire en vigueur} (prevailing wage) est le salaire que les règles du Department fixent comme minimum pour la profession, la zone et le niveau~;
l'employeur doit verser le plus élevé entre celui-ci et le salaire réel qu'il verse à des salariés comparables.
\end{definition}

\begin{itemize}
\item \textbf{Les labour condition applications} donnent le salaire annuel de base offert pour chaque poste parrainé. Un bonus compte
dans le salaire requis seulement s'il est «~assuré~», et non lorsqu'il est «~subordonné à un événement, comme les bénéfices annuels de
l'employeur~»~; le salaire proposé est donc le \omterm{def:in:how-pay-works:base}{salaire de base}. La population est celle des postes parrainés, et non de l'ensemble du personnel.
\item \textbf{L'enquête professionnelle} (OEWS) donne les centiles de salaire par secteur et par profession. Ses salaires incluent
«~les rémunérations incitatives, y compris les commissions et les bonus de production~» mais excluent les «~bonus hors production~», où se range le
bonus discrétionnaire, et elle marque les cellules les plus élevées comme plafonnées (top-coded).
\item \textbf{Les comptes déposés} donnent les \omterm{def:in:reading-a-trading-firms-accounts:staff}{charges de personnel} rapportées à l'\omterm{def:in:reading-a-trading-firms-accounts:staff}{effectif moyen} (chapitre~11)~: tous les salariés, chaque élément de
rémunération que les comptes de l'exercice imputent, et les charges sociales de l'employeur, pour une seule entité.
\item \textbf{Les déclarations de ratio de rémunération} des sociétés américaines cotées donnent la \omterm{def:in:how-pay-works:base}{rémunération totale} de leur salarié médian.
\item \textbf{Les déclarations de rémunération} des banques et des entreprises d'investissement donnent les effectifs et les totaux pour les mieux payés~:
les tableaux du Pilier~3 du chapitre~11 et les rapports sur les \omterm{def:in:pay-levels-by-role-firm-type-and-seniority:median}{hauts revenus} de l'Autorité bancaire européenne.
\end{itemize}

\begin{definition}[plafonnement des valeurs (top-coding), suppression des petites cellules (small-cell suppression)]\label{def:in:pay-levels-by-role-firm-type-and-seniority:suppress}
Le \emph{plafonnement des valeurs}\index{plafonnement des valeurs} (top-coding) est le remplacement des valeurs au-delà d'un seuil par un marqueur, de sorte que les cellules les plus élevées d'une enquête
n'indiquent que leur dépassement de ce seuil. La \emph{suppression des petites cellules}\index{suppression des petites cellules} (small-cell suppression) est la règle qui interdit de publier une statistique calculée à partir de trop peu d'observations,
parce qu'elle décrirait des individus identifiables ou serait trop bruitée
pour être utilisée.
\end{definition}

Les descriptions de champs de l'enquête professionnelle définissent encore \# comme «~un salaire égal ou supérieur à 115.00~\$ par heure ou
239\,200~\$ par an~», bien que le fichier de mai 2025 publie des valeurs au-delà de ce seuil (aucune des cellules utilisées dans ce livre n'est
marquée)~; ce chapitre applique sa propre suppression aux dossiers de visa~: aucune cellule n'est publiée à partir de moins de dix demandes, ou de
moins de trois employeurs, car une cellule tirée d'un ou deux employeurs publierait la rémunération d'une entreprise.

\begin{definition}[rémunération du salarié médian (median-employee pay), haut revenu (high earner)]\label{def:in:pay-levels-by-role-firm-type-and-seniority:median}
La \emph{rémunération du salarié médian}\index{rémunération du salarié médian} (median-employee pay) est la \omterm{def:in:how-pay-works:base}{rémunération totale} annuelle du salarié médian
d'une société américaine cotée, que sa déclaration de procuration publie avec le ratio de rémunération du directeur général.
Un \emph{haut revenu}\index{haut revenu} (high earner), au sens des statistiques bancaires européennes, est une personne payée un million d'euros ou plus au cours d'un exercice financier~;
les banques et les entreprises d'investissement transmettent leurs chiffres et leurs rémunérations à leurs superviseurs, et l'Autorité bancaire européenne
en publie les agrégats.
\end{definition}

\section{Construire des fourchettes à partir des dossiers de conditions de travail}

Le Department of Labor publie chaque demande qu'il tranche, trimestre après trimestre, dans des fichiers pouvant compter jusqu'à quelques centaines
de milliers de lignes. Le chapitre lit les quatre fichiers trimestriels des exercices 2021 et 2025, dédoublonnés par numéro de dossier.

\begin{definition}[Standard Occupational Classification, niveau de salaire (wage level)]\label{def:in:pay-levels-by-role-firm-type-and-seniority:soc}
La \emph{Standard Occupational Classification}\index{Standard Occupational Classification} (SOC) est la norme statistique fédérale américaine qui affecte chaque travailleur
à l'une des 867 professions détaillées~; 13-2099.01 désigne les «~Financial Quantitative Analysts~». Un
\emph{niveau de salaire}\index{niveau de salaire} (wage level) est l'un des quatre niveaux de \omterm{def:in:pay-levels-by-role-firm-type-and-seniority:lca}{salaire en vigueur} auxquels une demande est rattachée~:
I (débutant), II (qualifié), III (expérimenté) et IV (pleinement compétent), choisi à partir des exigences du poste.
\end{definition}

\begin{method}[Fourchettes à partir des labour condition applications]\label{met:in:pay-levels-by-role-firm-type-and-seniority:lca}
\begin{enumerate}
\item Diffuser chaque fichier ligne par ligne, en ne lisant que les colonnes nécessaires~; ne jamais lire les champs de contact.
\item Garder les demandes certifiées à temps plein.
\item Classer l'employeur d'après son nom, à l'aide de règles dont le modèle économique est sourcé dans un chapitre antérieur (douze teneurs de marché, sept fonds systématiques, quatre plateformes,
neuf banques, une société de gestion, quatre bourses), et le métier d'après l'intitulé du poste,
ou d'après le code SOC lorsque l'intitulé est générique.
\item Ramener la borne inférieure du salaire proposé à l'année, en fonction de son unité (année, mois, deux semaines, semaine, heure).
\item Regrouper par exercice fiscal, type d'employeur, métier et \omterm{def:in:pay-levels-by-role-firm-type-and-seniority:soc}{niveau de salaire}~; calculer les centiles et un intervalle bootstrap pour la
médiane~; supprimer les cellules de moins de dix demandes ou de moins de trois employeurs.
\end{enumerate}
\end{method}

\begin{dated}{2025-09}{Ce que contiennent les dossiers de visa pour les employeurs de ce livre}\label{dat:in:pay-levels-by-role-firm-type-and-seniority:counts}
\footnotesize
\begin{tabular}{@{}lrrrr@{}}
\toprule
type d'employeur & demandes exercice 2021 & employeurs & demandes exercice 2025 & employeurs\\
\midrule
banque & 3\,539 & 9 & 7\,298 & 9\\
teneur de marché & 259 & 10 & 367 & 10\\
fonds systématique & 143 & 5 & 210 & 7\\
plateforme multigérants & 53 & 3 & 102 & 3\\
bourse & 228 & 4 & 462 & 4\\
société de gestion & 200 & 1 & 369 & 1\\
\bottomrule
\end{tabular}

\smallskip\noindent Demandes certifiées à temps plein ayant un métier dans les familles du chapitre, exercices fiscaux américains (d'octobre à
septembre), issues des fichiers de divulgation du Department of Labor. Les cellules de la société de gestion proviennent d'un seul employeur et sont
toutes supprimées.
\end{dated}

Les banques déposent vingt fois plus de demandes que les teneurs de marché, parce qu'elles emploient beaucoup plus de personnes dans la
technologie~; une statistique regroupée sur l'ensemble du secteur serait une statistique bancaire. Chaque fourchette ci-dessous est donc ventilée par
type d'employeur.

\section{Fourchettes par métier, type d'entreprise et ancienneté}

\begin{center}\footnotesize
\begin{tabular}{@{}>{\raggedright\arraybackslash}p{4.0cm}>{\centering\arraybackslash}p{1.9cm}>{\centering\arraybackslash}p{1.9cm}>{\centering\arraybackslash}p{1.9cm}>{\centering\arraybackslash}p{1.9cm}>{\centering\arraybackslash}p{1.9cm}@{}}
\toprule
médiane du \omterm{def:in:how-pay-works:base}{salaire de base} offert, EF2025 (milliers de \$) & fonds systématique & plateforme & teneur de marché & banque & bourse\\
\midrule
tous les métiers & 215.0 & 190.5 & 175.0 & 154.3 & 130.8\\
chercheur quantitatif ou analyste & 220.0 & 190.0 & 175.0 & 158.1 & 107.1\\
ingénieur logiciel & 205.5 & 190.0 & 175.0 & 155.7 & 140.4\\
trader & -- & -- & 195.0 & 235.0 & --\\
\bottomrule
\end{tabular}

\smallskip Les cellules marquées -- sont supprimées. Données~: \texttt{data/industry/lca\_ranges.csv}.
\end{center}

L'ordre par type d'employeur est le même pour les deux métiers principaux~: les fonds systématiques offrent la base la plus élevée, puis les
plateformes, les teneurs de marché, les banques et les bourses, avec environ 84\,000~\$ d'écart entre le haut et le bas pour l'ensemble des métiers
réunis (\cref{fig:in:pay-levels-by-role-firm-type-and-seniority:kind}). Au sein d'un même type, la dispersion est forte~:
l'écart interquartile des chercheurs quantitatifs chez les teneurs de marché va de 132\,400~\$ à 237\,500~\$.
Entre l'exercice fiscal 2021 et 2025, la médiane tous métiers confondus a augmenté de 22.9~\% dans les fonds systématiques et dans les banques, de 18.7~\% dans les plateformes, de 16.7~\%
chez les teneurs de marché et de 7.9~\% dans les bourses~; les prix à la consommation aux États-Unis ont progressé de 18.8~\% sur les mêmes années civiles, si bien que, en termes
réels, la base offerte est restée stable ou a légèrement augmenté, et a baissé dans les bourses.

\begin{omfigure}
\begin{tikzpicture}
\begin{axis}[width=10.5cm,height=5.8cm,xmin=0.4,xmax=5.6,ymin=60,ymax=320,xtick=data,
  xticklabels={{fonds systématique},{plateforme},{teneur de marché},{banque},{bourse}},
  xticklabel style={font=\footnotesize,text width=1.8cm,align=center},yticklabel style={font=\footnotesize},
  ylabel={\small base offerte, milliers de \$},
  grid=major,grid style={gray!20},table/col sep=comma,unbounded coords=jump,
  legend style={legend cell align=left,font=\scriptsize,at={(0.5,-0.3)},anchor=north,/tikz/every even column/.append style={column sep=8pt}},
  legend columns=3]
\addplot+[only marks,mark=*,color=omDef,error bars/.cd,y dir=both,y explicit]
  table[x expr=\thisrow{pos}-0.18,y=all,y error plus=all_up,y error minus=all_dn]{figdata/industry/14-pay-levels-by-role-firm-type-and-seniority/base_by_kind.csv};
\addplot+[only marks,mark=square*,color=omThm,error bars/.cd,y dir=both,y explicit]
  table[x=pos,y=qr,y error plus=qr_up,y error minus=qr_dn]{figdata/industry/14-pay-levels-by-role-firm-type-and-seniority/base_by_kind.csv};
\addplot+[only marks,mark=triangle*,color=black!70,error bars/.cd,y dir=both,y explicit]
  table[x expr=\thisrow{pos}+0.18,y=se,y error plus=se_up,y error minus=se_dn]{figdata/industry/14-pay-levels-by-role-firm-type-and-seniority/base_by_kind.csv};
\legend{tous les métiers,chercheur quantitatif ou analyste,ingénieur logiciel}
\end{axis}
\end{tikzpicture}

\omcaption{Médiane du \omterm{def:in:how-pay-works:base}{salaire de base} annuel offert dans les demandes LCA américaines de l'exercice fiscal 2025, avec l'écart interquartile,
par type d'employeur et par métier. Base seule~: les bonus ne figurent pas dans ces dossiers. Données~: \texttt{data/industry/lca\_ranges.csv},
issues des fichiers de divulgation du Department of Labor, via \texttt{in\_lca\_derive.py}.}\label{fig:in:pay-levels-by-role-firm-type-and-seniority:kind}
\end{omfigure}

Le \omterm{def:in:pay-levels-by-role-firm-type-and-seniority:soc}{niveau de salaire} est la seule mesure d'ancienneté que contiennent les dossiers, et c'est une mesure faible là où la rémunération est très au-dessus du plancher
(\cref{fig:in:pay-levels-by-role-firm-type-and-seniority:level}). Dans les banques et les bourses, la médiane augmente avec le niveau,
comme le prévoient les définitions du Department~: dans les banques, de 100\,000~\$ au niveau~I à 169\,900~\$ au niveau~IV.
Chez les teneurs de marché, la médiane est de 84\,800~\$ au niveau~I et de 180\,000~\$ à 185\,000~\$ aux niveaux II à IV~: l'employeur
choisit le niveau qui décrit les exigences du poste, et lorsque l'offre se situe nettement au-dessus du \omterm{def:in:pay-levels-by-role-firm-type-and-seniority:lca}{salaire en vigueur} (prevailing wage) de chaque
niveau, ce choix ne contraint rien. Pour les chercheurs quantitatifs chez les teneurs de marché, la médiane du niveau~III
est inférieure à celle du niveau~II.

\begin{omfigure}
\begin{tikzpicture}
\begin{axis}[width=10.5cm,height=5.4cm,xmin=0.8,xmax=4.2,ymin=60,ymax=260,xtick={1,2,3,4},
  xticklabels={I,II,III,IV},tick label style={font=\footnotesize},
  xlabel={\small niveau de salaire~: I débutant, II qualifié, III expérimenté, IV pleinement compétent},
  ylabel={\small base médiane, milliers de \$},grid=major,grid style={gray!20},table/col sep=comma,
  unbounded coords=jump,legend style={legend cell align=left,font=\scriptsize,at={(0.5,-0.3)},anchor=north,
  /tikz/every even column/.append style={column sep=8pt}},legend columns=4]
\addplot[omDef,thick,mark=*] table[x=level,y=bank]{figdata/industry/14-pay-levels-by-role-firm-type-and-seniority/base_by_level.csv};
\addplot[omThm,thick,mark=square*] table[x=level,y=market_maker]{figdata/industry/14-pay-levels-by-role-firm-type-and-seniority/base_by_level.csv};
\addplot[black!70,thick,mark=triangle*] table[x=level,y=systematic_fund]{figdata/industry/14-pay-levels-by-role-firm-type-and-seniority/base_by_level.csv};
\addplot[gray,thick,dashed,mark=o,mark options={solid}] table[x=level,y=exchange]{figdata/industry/14-pay-levels-by-role-firm-type-and-seniority/base_by_level.csv};
\legend{banque,teneur de marché,fonds systématique,bourse}
\end{axis}
\end{tikzpicture}

\omcaption{Médiane du \omterm{def:in:how-pay-works:base}{salaire de base} offert selon le niveau de \omterm{def:in:pay-levels-by-role-firm-type-and-seniority:lca}{salaire en vigueur}, exercice fiscal 2025, tous métiers confondus. Le niveau suit la rémunération dans les banques et les
bourses, et s'aplatit au-dessus du niveau~I chez les teneurs de marché. La cellule du niveau~I des fonds systématiques est supprimée. Données~: comme
la \cref{fig:in:pay-levels-by-role-firm-type-and-seniority:kind}.}\label{fig:in:pay-levels-by-role-firm-type-and-seniority:level}
\end{omfigure}

L'enquête professionnelle est le contrôle extérieur au système des visas. En mai 2025, les développeurs de logiciels du secteur des valeurs
mobilières et des investissements avaient un salaire médian de 163\,290~\$ (écart interquartile de 132\,650~\$ à 203\,840~\$), et dans les banques
(intermédiation de crédit) 136\,620~\$~; les analystes financiers et d'investissement dans les valeurs mobilières 124\,370~\$. Les médianes des dossiers pour les
ingénieurs logiciels, 175\,000~\$ chez les teneurs de marché, 155\,700~\$ dans les banques et 140\,400~\$ dans les bourses, se situent dans la
même fourchette~: les postes parrainés de ces employeurs sont rémunérés comme ceux du secteur, et non en dessous.

\section{De la base au total~: ce que le coût du personnel par tête dit des bonus}

La base est la partie que montrent les dossiers~; la \omterm{def:in:how-pay-works:base}{rémunération totale} est celle qui intéresse le candidat. Le seul chiffre déclaré qui contienne
les bonus de l'ensemble d'une entreprise est le coût du personnel par tête, et le seul chiffre déclaré qui décrive un salarié typique est
la \omterm{def:in:pay-levels-by-role-firm-type-and-seniority:median}{rémunération du salarié médian}. Mis en regard l'un de l'autre et des dossiers de visa, ils encadrent le ratio manquant.

\begin{method}[Ratio implicite de la rémunération totale au salaire de base]\label{met:in:pay-levels-by-role-firm-type-and-seniority:ratio}
Pour un type d'employeur, on divise le coût du personnel par tête, tel qu'il figure dans les comptes déposés, par la médiane du \omterm{def:in:how-pay-works:base}{salaire de base} offert
de même type dans les dossiers de visa. Le résultat borne le ratio de la \omterm{def:in:how-pay-works:base}{rémunération totale} au \omterm{def:in:how-pay-works:base}{salaire de base} par le haut, car le coût du personnel inclut les
charges sociales de l'employeur, les avantages et l'acquisition des attributions différées, et c'est une moyenne sur une distribution asymétrique~; elle mélange
deux populations~: tous les salariés d'une seule entité face aux postes parrainés de plusieurs.
\end{method}

\begin{example}[Trois types d'employeurs]\label{ex:in:pay-levels-by-role-firm-type-and-seniority:ratio}
Le partenariat britannique d'un teneur de marché a dépensé entre 0.78~million de dollars et 2.11~millions de dollars par salarié en 2023--2025 (chapitre~11)
contre une médiane du \omterm{def:in:how-pay-works:base}{salaire de base} offert de 175\,000~\$ chez les teneurs de marché~: un ratio implicite de 4.4 à 12.0. Un \omterm{def:in:quantitative-hedge-funds-and-systematic-managers:sys}{gérant systématique} londonien
a dépensé 4.11~millions de dollars par salarié sur son exercice clos en janvier 2025, contre les 215\,000~\$ des fonds systématiques~: 19.1. Une
banque américaine, la rémunération et les avantages se situaient entre 342\,000~\$ et 399\,000~\$ par salarié en 2023--2025, contre les 154\,260~\$ des banques~:
2.2 à 2.6.
\end{example}

L'étendue du ratio selon les types d'entreprise est le résultat central du chapitre. Dans une banque, la \omterm{def:in:how-pay-works:base}{rémunération totale} n'est qu'un petit multiple de la base
pour le salarié moyen, parce que la banque emploie beaucoup de personnes dont la rémunération est surtout la base~; chez le teneur de marché et le
\omterm{def:in:quantitative-hedge-funds-and-systematic-managers:sys}{gérant systématique}, dont les effectifs sont réduits et dont le \omterm{def:in:revenue-and-profit-across-the-industry:perhead}{revenu par tête} est élevé (chapitre~12), la moyenne vaut plusieurs fois la base typique.
Les réserves jouent dans les deux sens~: les associés de l'entité sont hors de ses coûts de personnel, ce qui sous-estime le haut de l'échelle,
et sa moyenne inclut quelques très fortes attributions, ce qui surestime celle du salarié typique.

La \omterm{def:in:pay-levels-by-role-firm-type-and-seniority:median}{rémunération du salarié médian} montre à quel point la moyenne surestime le salarié typique, même là où la rémunération est modérée. Pour 2025, le
salarié médian d'Intercontinental Exchange percevait 112\,902~\$, contre un coût du personnel par tête de 152\,834~\$, soit 1.35 fois plus~;
chez Nasdaq, 101\,885~\$ contre 146\,142~\$ (1.43)~; chez Morgan Stanley, 136\,396~\$ contre 352\,000~\$ (2.58). Une moyenne sur l'ensemble des effectifs d'une
banque vaut deux fois et demie la rémunération de son salarié médian~; plus la rémunération d'une entreprise est concentrée,
plus l'écart est grand.

\section{La queue de la distribution~: nombre de hauts revenus et déclarations des preneurs de risque}

\begin{dated}{2026-04}{Les mieux rémunérés des banques et entreprises d'investissement de l'UE}\label{dat:in:pay-levels-by-role-firm-type-and-seniority:tail}
\footnotesize
\begin{tabular}{@{}>{\raggedright\arraybackslash}p{4.4cm}>{\centering\arraybackslash}p{2.6cm}>{\centering\arraybackslash}p{2.6cm}>{\centering\arraybackslash}p{2.6cm}@{}}
\toprule
2024, UE & \omterm{def:in:pay-levels-by-role-firm-type-and-seniority:median}{hauts revenus} & \omterm{def:in:how-pay-works:base}{rémunération totale} moyenne (\euro) & variable sur fixe\\
\midrule
banques, toutes activités & 2\,266 & 1\,830\,502 & 0.99\\
banques, banque d'investissement & 904 & 1\,757\,914 & 1.23\\
entreprises d'investissement, toutes activités & 288 & 2\,088\,931 & 3.59\\
entreprises d'investissement, négociation pour compte propre & 97 & 2\,094\,482 & 8.14\\
\bottomrule
\end{tabular}

\smallskip\noindent Autorité bancaire européenne, \omterm{def:in:pay-levels-by-role-firm-type-and-seniority:median}{hauts revenus} 2024, annexes II et III (publiées en avril 2026)~: 2\,554 \omterm{def:in:pay-levels-by-role-firm-type-and-seniority:median}{hauts
revenus} au total (2\,343 en 2023). Le rapport propre d'une banque britannique compte 717 de ses salariés payés 1~million d'euros ou plus en 2025
(chapitre~11).
\end{dated}

La queue de la distribution est l'endroit où les types d'entreprises diffèrent le plus. Parmi les \omterm{def:in:pay-levels-by-role-firm-type-and-seniority:median}{hauts revenus} de l'UE, ceux de la banque d'investissement des banques touchaient
en moyenne une part variable à peu près égale à la part fixe, parce que le plafonnement des bonus s'applique encore aux banques de l'UE (chapitre~15)~; ceux des
entreprises d'investissement pratiquant la négociation pour compte propre, que le plafonnement ne couvre pas, touchaient une part variable égale à 8.14 fois la part fixe, et
les 97 d'entre eux touchaient en moyenne 2.09~millions d'euros, soit environ 2.37~millions de dollars au taux moyen de 2025. Il s'agit de décomptes de personnes au-dessus d'un
seuil~: ils ne disent rien du nombre de personnes qui en sont proches, et ne couvrent que les établissements agréés dans l'UE.

\begin{omfigure}
\begin{tikzpicture}
\begin{axis}[width=10.5cm,height=5.6cm,xmode=log,xmin=100,xmax=4000,ymin=0.4,ymax=6.6,ytick=data,
  yticklabels={{banque, fixe proposé (Q1--Q3)},{teneur de marché, fixe proposé (Q1--Q3)},{banque, charges de personnel par tête},{teneur de marché, charges de personnel par tête},{UE, hauts revenus, banque d'investissement (moyenne)},{UE, hauts revenus, compte propre (moyenne)}},
  yticklabel style={font=\scriptsize},xticklabel style={font=\footnotesize},log ticks with fixed point,
  xtick={100,200,500,1000,2000},xlabel={\small milliers de \$ par an (échelle logarithmique)},grid=major,grid style={gray!20},
  table/col sep=comma]
\addplot+[only marks,mark=*,mark size=2pt,color=omDef,error bars/.cd,x dir=plus,x explicit,
  error bar style={line width=2.5pt,omDef!50},error mark=none]
  table[x=lo,y=pos,x error expr=\thisrow{hi}-\thisrow{lo}]{figdata/industry/14-pay-levels-by-role-firm-type-and-seniority/slices.csv};
\end{axis}
\end{tikzpicture}

\omcaption{Quatre tranches de rémunération sur une même échelle logarithmique~: le \omterm{def:in:how-pay-works:base}{salaire de base} proposé (écart interquartile, dossiers de visa, exercice 2025), les \omterm{def:in:reading-a-trading-firms-accounts:staff}{charges de personnel}
par tête issues des comptes déposés (fourchette sur 2023--2025~; une banque, un partenariat britannique de teneur de marché), et la rémunération moyenne
des \omterm{def:in:pay-levels-by-role-firm-type-and-seniority:median}{hauts revenus} de l'UE en 2024 au taux moyen dollar de 2025. Chacune mesure une chose différente sur une population différente.
Données~: \texttt{in\_paylevels} et les tableaux du chapitre.}\label{fig:in:pay-levels-by-role-firm-type-and-seniority:slices}
\end{omfigure}

\section{Tutoriel~: le pipeline des conditions de travail}\label{tut:in:pay-levels-by-role-firm-type-and-seniority:tut}

\begin{tutorial}
\textbf{Objectif.} Reconstruire les fourchettes du chapitre à partir des fichiers du Department of Labor, et les joindre aux autres tranches.
\textbf{État final~:} le tableau des fourchettes,
\cref{fig:in:pay-levels-by-role-firm-type-and-seniority:kind,fig:in:pay-levels-by-role-firm-type-and-seniority:level,fig:in:pay-levels-by-role-firm-type-and-seniority:slices}.
\begin{tutsteps}
\item \textbf{Récupérer.} Les huit fichiers trimestriels (exercices 2021 et 2025, environ 1.2~Go) sont placés dans un répertoire temporaire~; ils ne sont
jamais versionnés.
\item \textbf{Lire en flux.} \texttt{firm.paydata.read\_lca(path, keep)} produit une ligne à la fois à partir du classeur en lecture seule
(\cref{lst:in:pay-levels-by-role-firm-type-and-seniority:read})~; le script du chapitre
\texttt{in\_lca\_derive.py} met en cache les lignes retenues (type, rôle, niveau, salaire, rien d'autre) afin que les fichiers bruts ne soient lus qu'une fois,
sous une limite de mémoire.
\omcode{code/firm/paydata/firm_paydata.py}{36}{50}{Lire un classeur de déclarations ligne par ligne, en ne gardant que les
colonnes nécessaires.}\label{lst:in:pay-levels-by-role-firm-type-and-seniority:read}
\item \textbf{Classer et résumer.} \texttt{classify} avec \texttt{data/industry/lca\_employers.csv},
\texttt{classify\_role}, \texttt{annualise} et \texttt{cell} (centiles, intervalle bootstrap, suppression) écrivent
\texttt{lca\_ranges.csv}, \texttt{lca\_titles.csv} et \texttt{lca\_counts.csv}.
\item \textbf{Joindre.} \texttt{in\_paylevels.implied\_ratios} et \texttt{median\_ratios} placent les \omterm{def:in:reading-a-trading-firms-accounts:staff}{charges de personnel} des dépôts et
les médianes des déclarations de procuration face aux fourchettes. Les tests ne lisent que les tableaux versionnés et un classeur synthétique écrit
au moment du test.
\end{tutsteps}
\end{tutorial}

\textbf{Pour aller plus loin.} Ajouter les exercices 2022 à 2024 et suivre la médiane de chaque type d'entreprise à travers les hausses de taux~; ventiler les
ingénieurs logiciels des banques par État de travail et les comparer à l'enquête professionnelle de l'État.

\section{Construction~: le lecteur des données de rémunération}\label{bld:in:pay-levels-by-role-firm-type-and-seniority:paydata}

\begin{build}
\bfield{Rôle} Lire les sources publiques de rémunération pour en tirer des fourchettes et des enregistrements qui portent leurs réserves, pour ce chapitre, les chapitres consacrés
aux rôles (16--25) et le choix du chapitre~30.
\bfield{Interface} \texttt{firm.paydata}~: \texttt{LCA\_COLS}~; \texttt{read\_lca(path, keep)}~;
\texttt{annualise(amount, unit)}~; \texttt{Rule}, \texttt{load\_rules}, \texttt{classify(name, rules)}~; \texttt{ROLE\_RULES},
\texttt{classify\_role(title, soc)}~; \texttt{MIN\_CELL}~; \texttt{cell(values, rng, n\_boot)}~;
\texttt{PayEvidence(source, measure, population, year, lo, hi, unit, caveat)}. openpyxl en mode lecture seule et NumPy.
\bfield{Règles} Lire en flux, ne jamais charger un fichier entier~; ne lire aucun champ de contact ni d'information personnelle~; ne retenir que les demandes certifiées à temps plein~;
annualiser la borne basse de la fourchette proposée~; supprimer toute cellule de moins de dix demandes (le script du chapitre ajoute
trois employeurs)~; chaque source devient un enregistrement avec sa mesure, sa population et sa réserve.
\bfield{Tests d'acceptation} \texttt{code/firm/paydata/tests/}~: un classeur synthétique est lu en flux et filtré sans sa
colonne de contact~; les unités s'annualisent~; les règles d'employeur et de rôle classent et se replient sur le code SOC~; une cellule de neuf lignes est
supprimée et une cellule de cent lignes donne les bons centiles ainsi qu'un intervalle autour de sa médiane.
\bfield{Pour aller plus loin} L'enquête professionnelle par aire métropolitaine~; les dossiers PERM (certification permanente du travail), qui portent
des offres de postes élevés~; une estimation sur petite zone qui emprunte de la force aux différents types d'employeurs.
\end{build}

\begin{omsources}
\item US Department of Labor, OFLC, fichiers de déclaration LCA pour les exercices 2021 et 2025, et page des données de performance~;
20 CFR 655.731; OFLC Prevailing Wage Determination Policy Guidance (2009).
\item BLS, Standard Occupational Classification~; O*NET 13-2099.01~; OEWS mai 2025, estimations nationales par secteur
et notes techniques.
\item 17 CFR 229.402(u)~; déclarations de procuration pour 2025 d'Intercontinental Exchange, de Nasdaq, de Morgan Stanley et de Coinbase.
\item Autorité bancaire européenne, hauts revenus 2024 et annexes~; sa mention légale.
\item The Goldman Sachs Group, Form 10-K pour 2025~; chapitres 11 à 13 pour les comptes déposés et le rapport Pilier~3.
\end{omsources}

\section{Exercices}

\begin{exercise}[$\star$]\label{exo:in:pay-levels-by-role-firm-type-and-seniority:1}
Une demande indique un salaire de 95~\$ de l'heure. Quel \omterm{def:in:how-pay-works:base}{salaire de base} annuel la règle du chapitre enregistre-t-elle~? Et pour 16\,000~\$ par mois~?
\end{exercise}

\begin{exercise}[$\star$]\label{exo:in:pay-levels-by-role-firm-type-and-seniority:2}
En utilisant le tableau des médianes, de combien le \omterm{def:in:how-pay-works:base}{salaire de base} médian proposé pour tous les rôles des fonds systématiques dépasse-t-il
celui des bourses, en dollars et en pourcentage~?
\end{exercise}

\begin{exercise}[$\star$]\label{exo:in:pay-levels-by-role-firm-type-and-seniority:3}
Calculer le ratio entre les \omterm{def:in:reading-a-trading-firms-accounts:staff}{charges de personnel} par tête et la \omterm{def:in:pay-levels-by-role-firm-type-and-seniority:median}{rémunération du salarié médian} pour l'Intercontinental Exchange en 2025, à partir des \omterm{def:in:reading-a-trading-firms-accounts:staff}{charges de personnel}
de 1\,963 millions de dollars et de 12\,844 salariés.
\end{exercise}

\begin{exercise}[$\star\star$]\label{exo:in:pay-levels-by-role-firm-type-and-seniority:4}
Pourquoi un bonus ne peut-il pas faire partie du salaire dans une \omterm{def:in:pay-levels-by-role-firm-type-and-seniority:lca}{labor condition application}, et que mesurent alors les dépôts~?
\end{exercise}

\begin{exercise}[$\star\star$]\label{exo:in:pay-levels-by-role-firm-type-and-seniority:5}
La médiane du niveau III des teneurs de marché pour les chercheurs quantitatifs est inférieure à celle du niveau II. Donner deux explications
qui ne reposent pas sur le fait de payer moins les personnes expérimentées.
\end{exercise}

\begin{exercise}[$\star\star$]\label{exo:in:pay-levels-by-role-firm-type-and-seniority:6}
Déflater le \omterm{def:in:how-pay-works:base}{salaire de base} médian proposé des teneurs de marché de l'exercice fiscal 2021 aux dollars de 2025 avec l'indice des prix à la consommation (114.33
en 2021, 135.83 en 2025), puis comparer au salaire médian de l'exercice fiscal 2025.
\end{exercise}

\begin{exercise}[$\star\star\star$]\label{exo:in:pay-levels-by-role-firm-type-and-seniority:7}
\emph{Programmation.} À l'aide de \texttt{lca\_ranges.csv}, lister chaque cellule de l'exercice fiscal 2025 des teneurs de marché qui est supprimée, avec
son nombre de demandes et d'employeurs, et indiquer la règle qui a supprimé chacune.
\end{exercise}

\begin{exercise}[$\star\star\star$]\label{exo:in:pay-levels-by-role-firm-type-and-seniority:8}
\emph{Trouver la faille.} «~Le \omterm{def:in:how-pay-works:base}{salaire de base} médian des teneurs de marché est de 175\,000~\$ et leurs \omterm{def:in:reading-a-trading-firms-accounts:staff}{charges de personnel} par tête de 2.1 millions de dollars, donc un jeune diplômé
peut s'attendre à une \omterm{def:in:how-pay-works:base}{rémunération totale} d'environ 2 millions de dollars.~»
\end{exercise}

\section{Problème~: quatre sources, un seul nombre}

\begin{problem}[{Problème du week-end --- quatre sources, un seul nombre}]\label{pb:in:pay-levels-by-role-firm-type-and-seniority:1}
Un candidat qui hésite entre une banque et un teneur de marché veut un seul nombre~: de combien la \omterm{def:in:how-pay-works:base}{rémunération totale} dépasse le \omterm{def:in:how-pay-works:base}{salaire de base} chez chacun.
Aucune source ne le donne. Il faut le construire.

\medskip\noindent\textbf{Partie I --- Les sources.}
\begin{enumerate}
\item Définir une \omterm{def:in:pay-levels-by-role-firm-type-and-seniority:lca}{labor condition application} (LCA) et le \omterm{def:in:pay-levels-by-role-firm-type-and-seniority:lca}{salaire en vigueur} (prevailing wage). Que mesure le salaire figurant dans le dépôt~?
\item Que l'enquête sur les professions inclut-elle et exclut-elle des salaires, et que plafonne-t-elle~?
\item Définir le \omterm{def:in:pay-levels-by-role-firm-type-and-seniority:suppress}{plafonnement des valeurs} (top-coding) et la \omterm{def:in:pay-levels-by-role-firm-type-and-seniority:suppress}{suppression des petites cellules}, puis énoncer la règle de suppression du chapitre.
\item Définir la \omterm{def:in:pay-levels-by-role-firm-type-and-seniority:median}{rémunération du salarié médian} et un \omterm{def:in:pay-levels-by-role-firm-type-and-seniority:median}{haut revenu}.
\item Pour chacune des cinq sources, indiquer ce qu'elle mesure et sur qui.
\end{enumerate}

\medskip\noindent\textbf{Partie II --- Les fourchettes.}
\begin{enumerate}[resume]
\item Énoncer la méthode de construction des fourchettes à partir des dépôts.
\item Donner le nombre de demandes et d'employeurs par catégorie pour l'exercice fiscal 2025.
\item Donner le \omterm{def:in:how-pay-works:base}{salaire de base} médian proposé par catégorie pour tous les rôles, les chercheurs quantitatifs et les ingénieurs logiciels.
\item Comment les médianes ont-elles évolué de l'exercice fiscal 2021 à 2025, en valeur nominale et après inflation~?
\item Que montre le \omterm{def:in:pay-levels-by-role-firm-type-and-seniority:soc}{niveau de salaire} chez les banques et chez les teneurs de marché, et pourquoi cet écart~?
\end{enumerate}

\medskip\noindent\textbf{Partie III --- Du \omterm{def:in:how-pay-works:base}{salaire de base} au total.}
\begin{enumerate}[resume]
\item Énoncer la méthode du ratio implicite de la \omterm{def:in:how-pay-works:base}{rémunération totale} au \omterm{def:in:how-pay-works:base}{salaire de base}, ainsi que ses biais.
\item Le calculer pour le teneur de marché, le \omterm{def:in:quantitative-hedge-funds-and-systematic-managers:sys}{gérant systématique} et la banque.
\item Comparer les \omterm{def:in:reading-a-trading-firms-accounts:staff}{charges de personnel} par tête à la \omterm{def:in:pay-levels-by-role-firm-type-and-seniority:median}{rémunération du salarié médian} chez trois entreprises cotées.
\item Comment l'enquête sur les professions se compare-t-elle aux dépôts pour les ingénieurs logiciels~?
\item Qu'apportent les chiffres des \omterm{def:in:pay-levels-by-role-firm-type-and-seniority:median}{hauts revenus} dans l'UE, et que ne peuvent-ils pas dire~?
\end{enumerate}

\medskip\noindent\textbf{Partie IV --- Le verdict.}
\begin{enumerate}[resume]
\item Énoncer le \emph{résultat nommé}~: le ratio implicite de la \omterm{def:in:how-pay-works:base}{rémunération totale} au \omterm{def:in:how-pay-works:base}{salaire de base} par type d'employeur, sous forme de fourchette avec ses réserves.
\item Quelle réserve gonfle le plus le ratio du teneur de marché pour un salarié typique~?
\item Quelle tranche réduirait le plus la fourchette si elle était publique~?
\item Que doit demander le candidat à chaque employeur pour remplacer le ratio par des faits~?
\item En deux phrases, répondre au candidat.
\end{enumerate}
\end{problem}

\section{Questions d'entretien}

\begin{interviewq}[$\star$ \iqroles{chercheur, ingénieur ML}]\label{iq:in:pay-levels-by-role-firm-type-and-seniority:1}
Vous disposez de salaires exprimés en cinq unités (année, mois, deux semaines, semaine, heure). Écrivez la conversion et dites ce que vous feriez
d'une fourchette «~de 150\,000~\$ à 250\,000~\$~».
\end{interviewq}

\begin{interviewq}[$\star$ \iqroles{chercheur}]\label{iq:in:pay-levels-by-role-firm-type-and-seniority:2}
Pourquoi présenter un intervalle bootstrap pour une médiane plutôt qu'une erreur type~?
\end{interviewq}

\begin{interviewq}[$\star\star$ \iqroles{développeur}]\label{iq:in:pay-levels-by-role-firm-type-and-seniority:3}
Concevez un job qui lit un tableur de 300~MB en mémoire bornée et n'écrit que des agrégats. Qu'est-ce qui peut mal tourner~?
\end{interviewq}

\begin{interviewq}[$\star\star$ \iqroles{chercheur, risque}]\label{iq:in:pay-levels-by-role-firm-type-and-seniority:4}
Une cellule compte 12 observations provenant de deux employeurs. Devriez-vous publier sa médiane~? Pourquoi ou pourquoi pas~?
\end{interviewq}

\begin{interviewq}[$\star\star$ \iqroles{banque, trader}]\label{iq:in:pay-levels-by-role-firm-type-and-seniority:5}
Pourquoi la part variable peut-elle être à peu près égale à la part fixe pour les banquiers d'investissement d'une banque dans l'UE,
mais huit fois la part fixe dans les entreprises d'investissement négociant pour compte propre~?
\end{interviewq}

\begin{interviewq}[$\star\star\star$ \iqroles{chercheur}]\label{iq:in:pay-levels-by-role-firm-type-and-seniority:6}
Les dépôts de visa ne couvrent que les postes parrainés. Comment la sélection pourrait-elle biaiser leurs fourchettes de salaire, et comment le testeriez-vous~?
\end{interviewq}
